\documentclass[11pt,a4paper,fontset=fandol]{ctexart}

\usepackage[a4paper,top=2.4cm,bottom=2.6cm,left=2.35cm,right=2.35cm,headheight=18pt,headsep=0.55cm,footskip=26pt]{geometry}
\usepackage{amsmath,amssymb}
\usepackage{fancyhdr}
\usepackage{lastpage}
\usepackage{enumitem}
\usepackage{titlesec}
\usepackage{xcolor}
\usepackage{hyperref}
\usepackage{bookmark}
\usepackage[most]{tcolorbox}
\usepackage{setspace}
\usepackage{array}

\hypersetup{
  colorlinks=true,
  linkcolor=blue!50!black,
  urlcolor=blue!50!black
}

\setstretch{1.16}
\setlength{\parindent}{2em}
\setlength{\parskip}{0.32em}
\setlist[enumerate]{itemsep=0.35em, topsep=0.35em, leftmargin=2.3em}
\setlist[itemize]{itemsep=0.25em, topsep=0.25em, leftmargin=2em}

\definecolor{mainblue}{RGB}{36,83,140}
\definecolor{lightblue}{RGB}{241,246,252}
\definecolor{lightgray}{RGB}{246,246,246}
\definecolor{rulegray}{RGB}{110,118,128}
\definecolor{softgreen}{RGB}{236,247,240}
\definecolor{softyellow}{RGB}{254,249,229}

\titleformat{\section}
  {\Large\bfseries\color{mainblue}}
  {\thesection.}
  {0.45em}
  {}
  [\vspace{0.25em}\color{rulegray}\titlerule]
\titlespacing*{\section}{0pt}{1.2em}{0.8em}
\titleformat{\subsection}{\large\bfseries\color{mainblue}}{\thesubsection}{0.5em}{}

\pagestyle{fancy}
\fancyhf{}
\fancyhead[L]{\small 抽屉原理与极值思想专题目录页}
\fancyhead[C]{\small 组合专题资料索引}
\fancyhead[R]{\small 刘文博}
\fancyfoot[C]{\small 第 \thepage 页 / 共 \pageref{LastPage} 页}
\renewcommand{\headrulewidth}{0.55pt}
\renewcommand{\footrulewidth}{0.35pt}

\newtcolorbox{goalbox}[1]{
  breakable,
  colback=lightblue,
  colframe=mainblue,
  boxrule=0.7pt,
  arc=1mm,
  title=\textbf{#1},
  fonttitle=\bfseries
}

\newtcolorbox{infobox}[1]{
  breakable,
  colback=softgreen,
  colframe=green!40!black,
  boxrule=0.65pt,
  arc=1mm,
  title=\textbf{#1},
  fonttitle=\bfseries
}

\newtcolorbox{warnbox}[1]{
  breakable,
  colback=softyellow,
  colframe=orange!60!black,
  boxrule=0.65pt,
  arc=1mm,
  title=\textbf{#1},
  fonttitle=\bfseries
}

\begin{document}

\thispagestyle{empty}
\begin{titlepage}
\centering
\vspace*{0.9cm}

\begin{tcolorbox}[
  enhanced,
  width=0.92\textwidth,
  colback=white,
  colframe=mainblue,
  boxrule=1pt,
  arc=0mm,
  left=6mm,
  right=6mm,
  top=8mm,
  bottom=8mm
]
\centering

{\fontsize{24}{30}\selectfont\bfseries 抽屉原理与极值思想专题目录页\par}
\vspace{0.65cm}
{\fontsize{17}{22}\selectfont 讲义、题单与周测的使用说明\par}

\vspace{1cm}
\begin{tcolorbox}[colback=lightblue,colframe=mainblue,width=0.84\textwidth,boxrule=1pt]
\centering
{\large 适合对象：已经掌握基础组合与数论分类，准备系统学习“最少保证”和“存在性证明”的学生}\\[0.35em]
{\normalsize 本目录页帮助把抽屉原理与极值思想专题资料串成一个完整训练包}
\end{tcolorbox}

\vspace{1.0cm}
\renewcommand{\arraystretch}{1.42}
\begin{tabular}{>{\bfseries}p{3.5cm}p{8cm}}
专题名称： & 抽屉原理与极值思想 \\
资料组成： & 课堂讲义、提高训练题单、周测 A 卷、周测 B 卷 \\
使用方式： & 先学讲义，再做题单，最后用 A/B 周测检查掌握情况 \\
编译方式： & XeLaTeX \\
\end{tabular}

\vspace{0.9cm}
{\large \today\par}
\end{tcolorbox}
\end{titlepage}

\tableofcontents
\newpage

\section{专题包包含哪些资料}

\begin{goalbox}{四份核心资料}
\begin{enumerate}
  \item 课堂讲义：帮助理解基础抽屉原理、加强版抽屉、最少保证和极值法入口；
  \item 提高训练题单：用于课后巩固和变式提升，重点训练“会找抽屉、会估上界、会做经典存在性证明”；
  \item 周测 A 卷：检查基础分组和最少保证是否稳定；
  \item 周测 B 卷：检查综合模型、极值分析和经典构造是否真正掌握。
\end{enumerate}
\end{goalbox}

\begin{infobox}{对应文件名}
\begin{itemize}
  \item 讲义：\texttt{pigeonhole\_extremal\_handout.tex} / \texttt{pigeonhole\_extremal\_handout.pdf}
  \item 提高训练题单：\texttt{pigeonhole\_extremal\_advanced\_problem\_set.tex} / \texttt{pigeonhole\_extremal\_advanced\_problem\_set.pdf}
  \item 周测 A 卷：\texttt{pigeonhole\_extremal\_weekly\_test\_A.tex} / \texttt{pigeonhole\_extremal\_weekly\_test\_A.pdf}
  \item 周测 B 卷：\texttt{pigeonhole\_extremal\_weekly\_test\_B.tex} / \texttt{pigeonhole\_extremal\_weekly\_test\_B.pdf}
\end{itemize}
\end{infobox}

\section{建议学习顺序}

\begin{goalbox}{推荐顺序}
\begin{enumerate}
  \item 先完整学习讲义，重点吃透“物体是什么、抽屉是什么、最大安全数怎么求”；
  \item 再完成提高训练题单，训练余数分组、区间分组、几何分割和极值法；
  \item 接着做周测 A 卷，检查基础模型是否稳定；
  \item 最后做周测 B 卷，检查经典构造题和较难证明题是否掌握。
\end{enumerate}
\end{goalbox}

\begin{warnbox}{不建议的做法}
\begin{itemize}
  \item 一上来只背“抽屉原理”四个字，却不先想抽屉怎么构造；
  \item 直接写结论，不说明“安全状态最多能到哪里”；
  \item 极值题只会说“取最大值”，不会继续利用它推出新的限制；
  \item 几何型抽屉题不先分格子，而是盲目画点尝试。
\end{itemize}
\end{warnbox}

\section{每份资料主要解决什么问题}

\subsection{课堂讲义}
讲义适合第一次系统学习这个专题，主要解决以下问题：
\begin{itemize}
  \item 怎样从题目中识别“谁是物体、谁是抽屉”；
  \item 怎样处理“至少多少个才能保证出现某种现象”；
  \item 加强版抽屉原理和平均数思想为什么本质一致；
  \item 极值思想在组合题里怎样和抽屉法配合使用。
\end{itemize}

\subsection{提高训练题单}
题单适合第二阶段训练，主要用来把方法做熟。它更强调：
\begin{itemize}
  \item 能否按余数、区间、颜色或链分解来构造抽屉；
  \item 能否正确求出“最大安全数”；
  \item 能否处理几何型抽屉和线段分割问题；
  \item 能否用极值法给出严谨的上界证明。
\end{itemize}

\subsection{周测 A 卷}
A 卷偏基础，适合检查：
\begin{itemize}
  \item 余数分类是否熟练；
  \item 最少保证题是否会做；
  \item 线段、盒子、连续整数这些基础模型是否稳定；
  \item 极值题中的相邻差估计是否会写。
\end{itemize}

\subsection{周测 B 卷}
B 卷偏综合，适合检查：
\begin{itemize}
  \item 经典抽屉模型是否会推广到一般 $n$；
  \item 链分解和倍数关系是否会处理；
  \item 证明题是否真正解释清楚了“为什么这样分组”；
  \item 上界得出后，是否还能给出达到上界的构造。
\end{itemize}

\section{一个推荐的 4 次训练安排}

\begin{goalbox}{4 次完成一个小专题}
\begin{enumerate}
  \item 第 1 次：学习讲义前半部分，重点理解基础抽屉和最少保证；
  \item 第 2 次：学习讲义后半部分，完成题单基础题和中档题；
  \item 第 3 次：完成题单综合题，并把错题按“抽屉找错/上界错/极值关系漏写/构造不完整”分类订正；
  \item 第 4 次：完成周测 A 卷或 B 卷，作为阶段性检查。
\end{enumerate}
\end{goalbox}

\begin{infobox}{如果孩子基础较好，可以这样压缩}
\begin{itemize}
  \item 第 1 天：讲义 + 题单基础题；
  \item 第 2 天：题单变式题 + 错题订正；
  \item 第 3 天：周测 A 卷；
  \item 第 4 天：周测 B 卷。
\end{itemize}
\end{infobox}

\section{专题学习后的自检清单}

\begin{goalbox}{如果下面 8 条都能做到，说明这个专题基本过关}
\begin{enumerate}
  \item 我会先判断谁是物体、谁是抽屉；
  \item 我会处理最少保证题；
  \item 我会把数按余数类、区间或链来分类；
  \item 我会使用加强版抽屉原理；
  \item 我会处理线段分割和几何型抽屉题；
  \item 我会把“最多能选多少”写成严谨的上界证明；
  \item 我会给出达到上界的构造例子；
  \item 我会在解答后检查分类是否完整、是否真的证明了“保证”。
\end{enumerate}
\end{goalbox}

\section{给家长的使用建议}

\begin{warnbox}{更有效的带练方式}
\begin{itemize}
  \item 不要一开始就说“这是抽屉原理”，先让孩子自己说抽屉怎么分；
  \item 如果卡住，可以只提示“先想最安全的情况最多能取几个”；
  \item 极值题订正时，重点看孩子有没有把相邻差累加成总差；
  \item 做完后请孩子口头复述：这题到底是按余数分、按区间分，还是按链分解解决的。
\end{itemize}
\end{warnbox}

\begin{infobox}{和后续专题的衔接}
这个专题学完以后，最自然的下一步就是继续学习\textbf{染色方法与棋盘问题}。因为颜色本身也可以看成一种分类，而覆盖题和路径题里很多结论都可以看成抽屉、不变量与分类思想的进一步发展。
\end{infobox}

\end{document}
